% META: {"id": "TCOMPL-MF-CO-010", "chapitre": "TCOMPL-MODELES-FONCTION", "type_objet": "corrige", "exercice_ref": "TCOMPL-MF-EX-010", "capacites_codes": ["C7"], "sources_inspiration": [], "status": "generated"}
% BEGIN-VERIFY
% from sympy import Rational, simplify
% xs = [1, 2, 3, 4, 5]
% ys = [12, 15, 19, 22, 27]
% n = len(xs)
% xbar = Rational(sum(xs), n)
% ybar = Rational(sum(ys), n)
% assert xbar == 3 and ybar == 19
% cov = Rational(sum((x - xbar) * (y - ybar) for x, y in zip(xs, ys)), n)
% var = Rational(sum((x - xbar) ** 2 for x in xs), n)
% assert var == 2
% a = cov / var
% b = ybar - a * xbar
% assert a == Rational(37, 10)
% assert b == Rational(79, 10)
% assert simplify(a * 7 + b - Rational(338, 10)) == 0
% END-VERIFY
\begin{corrige}{TCOMPL-MF-EX-010}
\textbf{1.} $\bar x=3$ et $\bar y=19$ : $G(3\,;\,19)$.

\textbf{2.} La covariance vaut
$\tfrac15\sum(x_i-\bar x)(y_i-\bar y)=\tfrac{37}{10}$ et la variance de $x$
vaut $2$. La pente est $a=\tfrac{37}{10\times2}=\tfrac{37}{10}$ ; en fait, le
calcul exact donne $a=3{,}7$ et $b=\bar y-a\bar x=19-11{,}1=7{,}9$. La droite
est $y=3{,}7x+7{,}9$.

\textbf{3.} Au rang $7$ : $y=3{,}7\times7+7{,}9=33{,}8$, soit environ
$33\,800$ abonnés. L'extrapolation reste proche de la plage observée
($x\leq5$), mais rien ne garantit que la croissance linéaire se poursuive :
la prévision est indicative, pas certaine.
\end{corrige}
